% GENERATED by scripts/make_tables.py from config/implementation-register.toml.
% Do not edit by hand.
\begin{landscape}
\begingroup
\footnotesize
\setlength{\tabcolsep}{4pt}
\renewcommand{\arraystretch}{1.25}
\begin{longtable}{@{}p{24mm}p{13mm}p{34mm}p{30mm}p{21mm}p{30mm}p{28mm}p{28mm}@{}}
\caption{The couplings of this reconstruction written in each external convention. Every cell is derived from the line of \texttt{provenance/conventions.md} named in the \emph{Sheet} column; no cell is blank, and a treatment with no established counterpart says so rather than being left to inference. \textbf{The thesis takes two columns and they differ on exactly the kinematic rows} --- the $\dot{\mathcal T}$ form moves every kinematic term to the right-hand side, flipping its overall sign against the K/$\Pi$ form. On the remaining rows the two agree, and the table says \emph{as Th-K} rather than restating them: a difference in wording where there is none in substance is how two conventions get conflated while appearing distinct.}\label{tab:translations}\\
\toprule
\textbf{Coupling} & \textbf{Sheet} & \textbf{MGE99} & \textbf{CL} & \textbf{Th, K/$\Pi$} & \textbf{Th, $\dot{\mathcal T}$} & \textbf{BF} & \textbf{P09} \\
\midrule
\endfirsthead
\toprule
\textbf{Coupling} & \textbf{Sheet} & \textbf{MGE99} & \textbf{CL} & \textbf{Th, K/$\Pi$} & \textbf{Th, $\dot{\mathcal T}$} & \textbf{BF} & \textbf{P09} \\
\midrule
\endhead
\bottomrule
\endfoot
Rank raising & \texttt{C4.2} & identity & identity in form; $\mathrm D_a$ unchanged, the sign sits in $k$ & identity & as Th-K & not established: BF work in $(\ell,m)$ tetrad form, no coefficient reproduced & not established: Pitrou works in Poisson gauge, no hierarchy coefficient checked \\
Rank lowering & \texttt{C4.4} & identity & identity in form; CL (2.25) prints $+\mathrm D^bI_{bA_\ell}$, reproduced from MGE99 (71) & identity; Th App.~D footnote to (D.44) at $\ell=1$ agrees & as Th-K & not established & not established \\
Free-streaming mode recursion & \texttt{C4.5} & identity & identity in form; CL (2.25) free streaming reproduced from MGE99 (71) & identity & as Th-K & not established & not established \\
PSTF contraction on the sphere & \texttt{C6.1} & identity (one source: EMT 1983, via C6.4) & identity; CL expands in the same $O^{A_\ell}$ & identity, Th (1.20) & as Th-K & not established: BF use ${}_sY_{\ell m}$, not a PSTF basis & identity at the basis level: Pitrou (1.28) expands in $n^{A_\ell}$ \\
PSTF self-contraction $\beta_\ell$ & \texttt{C5.2} & identity & identity; CL absorbs it into $\Delta^C_\ell$, see the row below & identity, Th (1.23) & as Th-K & no counterpart: $(\ell,m)$ normalisation carries $\sqrt{4\pi/(2\ell+1)}$ instead & identity, Pitrou (1.29) in the equivalent $4\pi\ell!/(2\ell+1)!!$ form \\
Angular measure $\Delta_\ell$ & \texttt{C5.5} & identity, MGE99 (47), same value & $\Delta^C_\ell=4\pi(-2)^\ell(\ell!)^2/(2\ell+1)!=(-1)^\ell\Delta_\ell$, CL (2.23) & identity & as Th-K & no counterpart in this form & identity, Pitrou (1.29), verified symbolically \\
Shear coupling at $\ell+2$ & \texttt{C9, T-MGE and T-Th-T} & $-(\ell+2)$ as printed in (71), (88), (90), (99) --- \textbf{defect E9}; $+(\ell+2)$ is correct, and MGE99 (89) at $\ell=2$ prints $+\tfrac{8}{15}\rho\sigma$, agreeing with the correction & $-\tfrac{8}{15}\delta^2_\ell I\sigma$ gives $+(\ell+2)$, independently of MGE99 & $-(\ell+2)$ in Th (2.73) --- carries E9 & $+(\ell+2)$ \textbf{on the RHS} in Th (2.84): every kinematic term flips overall sign against (2.73) because it is moved across. \textbf{This row is where the two thesis columns differ} & not established: BF put the observer at rest, so they carry no shear coupling & not established \\
Acceleration coupling at $\ell\pm1$ & \texttt{C9, T-CL and T-Th-T} & identity & $A_a\to-A_a$ under $g\to-g$ at fixed $u^a$ (CL footnote to \S2), the signature difference and nothing more & identity & overall sign flipped with the rest of the kinematic terms --- \textbf{differs from Th-K} & not established: observer at rest & not established \\
\end{longtable}
\endgroup
\end{landscape}
